%% Đề 38 -- de-tuyen-sinh-lop-10-mon-toan-nam-2025-2026-so-gddt-lang-son.pdf
\deheader{ĐỀ SỐ 38}%
  {Kỳ thi Tuyển sinh vào lớp 10 THPT — Năm học 2025--2026}%
  {Sở GD\&ĐT Lạng Sơn | Chỉ phần Đại số}
\addcontentsline{toc}{section}{Đề 38 — Tuyển sinh vào lớp 10 THPT — Năm học 2025--2026}

\begin{de_cau}{Câu 1}{2,0 điểm}
\begin{enumerate}[label=\arabic*.]
\item Tính giá trị của biểu thức:
\[A = \sqrt{100} - \sqrt{64};\qquad B = \sqrt{(3-\sqrt{3})^2} + \sqrt{3}.\]
\item Cho biểu thức: $Q = \left(\dfrac{\sqrt{x}-3}{\sqrt{x}+2} + \dfrac{1}{\sqrt{x}-2}\right) : \dfrac{1}{x-4}$ với $x \ge 0,\; x \neq 4$.
\begin{enumerate}[label=\alph*)]
\item Rút gọn biểu thức $Q$.
\item Tìm $x$ để $Q = \dfrac{4}{5}$.
\end{enumerate}
\end{enumerate}
\end{de_cau}

\begin{loi_giai_cau}
\textbf{1.}
\[A = \sqrt{100} - \sqrt{64} = 10 - 8 = \boxed{2}.\]
\[B = |3-\sqrt{3}| + \sqrt{3} = 3-\sqrt{3}+\sqrt{3} = \boxed{3}.\]

\textbf{2.} Với $x \ge 0, x \neq 4$:
\[Q = \frac{(\sqrt{x}-3)(\sqrt{x}-2) + (\sqrt{x}+2)}{x-4} \cdot (x-4) = (x - 5\sqrt{x} + 6 + \sqrt{x} + 2) = \boxed{x - 4\sqrt{x} + 8}.\]
\textbf{b)} Để $Q = \dfrac{4}{5}$, giải phương trình bậc hai ẩn $\sqrt{x}$ và đối chiếu ĐKXĐ để tìm nghiệm.
\end{loi_giai_cau}

\begin{de_cau}{Câu 2}{2,0 điểm}
\begin{enumerate}[label=\arabic*.]
\item Giải hệ phương trình: 
\[ \begin{cases} 3x + 2y = 7 \\ x - 2y = 1 \end{cases} \]
\item Cho hàm số $y = x^2$ có đồ thị là $(P)$ và đường thẳng $(d)\colon y = 3x - 2$. Tìm tọa độ giao điểm của $(P)$ và $(d)$ bằng phép tính.
\end{enumerate}
\end{de_cau}

\begin{loi_giai_cau}
\textbf{1.} Cộng hai phương trình: $4x = 8 \implies x = 2$.
Thay vào phương trình thứ hai: $2 - 2y = 1 \implies 2y = 1 \implies y = \dfrac{1}{2}$.
Vậy nghiệm của hệ là $\boxed{\left(2; \dfrac{1}{2}\right)}$.

\textbf{2.} Phương trình hoành độ giao điểm:
\[x^2 = 3x - 2 \iff x^2 - 3x + 2 = 0 \iff (x-1)(x-2) = 0 \iff x = 1\text{ hoặc }x = 2.\]
\begin{itemize}
\item $x = 1 \implies y = 1 \implies A(1; 1)$.
\item $x = 2 \implies y = 4 \implies B(2; 4)$.
\end{itemize}
Tọa độ giao điểm là $\boxed{A(1; 1)}$ và $\boxed{B(2; 4)}$.
\end{loi_giai_cau}

\begin{de_cau}{Câu 3}{1,5 điểm}
Hai tổ công nhân cùng làm chung một công việc thì sau 12 giờ hoàn thành. Nếu tổ I làm trong 4 giờ rồi tổ II làm tiếp trong 7 giờ thì cả hai tổ hoàn thành được $50\%$ công việc. Hỏi nếu làm riêng thì mỗi tổ hoàn thành công việc trong bao lâu?
\end{de_cau}

\begin{loi_giai_cau}
Gọi thời gian tổ I làm một mình xong công việc là $x$ (giờ), tổ II làm một mình xong là $y$ (giờ) ($x, y > 12$).
Trong 1 giờ:
\begin{itemize}
\item Tổ I làm được $\dfrac{1}{x}$ công việc.
\item Tổ II làm được $\dfrac{1}{y}$ công việc.
\end{itemize}
Vì hai tổ làm chung trong 12 giờ thì xong nên: $\dfrac{1}{x} + \dfrac{1}{y} = \dfrac{1}{12} \quad (1)$.
Tổ I làm 4 giờ, tổ II làm 7 giờ được $50\% = \dfrac{1}{2}$ công việc:
\[\frac{4}{x} + \frac{7}{y} = \frac{1}{2} \quad (2).\]
Giải hệ phương trình bằng cách đặt $u = 1/x, v = 1/y$:
\[\begin{cases} u + v = 1/12 \\ 4u + 7v = 1/2 \end{cases} \implies \begin{cases} 4u + 4v = 1/3 \\ 4u + 7v = 1/2 \end{cases} \implies 3v = \frac{1}{6} \implies v = \frac{1}{18} \implies y = 18.\]
Suy ra $u = \dfrac{1}{12} - \dfrac{1}{18} = \dfrac{1}{36} \implies x = 36$.
Vậy tổ I làm riêng mất $\boxed{36\text{ giờ}}$, tổ II làm riêng mất $\boxed{18\text{ giờ}}$.
\end{loi_giai_cau}

\begin{de_cau}{Câu 6}{1,0 điểm}
Cho các số thực dương $a, b, c$ thỏa mãn $a + b + c = 3$. Tìm giá trị nhỏ nhất của:
\[P = \dfrac{a^2}{b+1} + \dfrac{b^2}{c+1} + \dfrac{c^2}{a+1}.\]
\end{de_cau}

\begin{loi_giai_cau}
Áp dụng bất đẳng thức Cauchy-Schwarz dạng Engel:
\[P \ge \frac{(a+b+c)^2}{(b+1) + (c+1) + (a+1)} = \frac{(a+b+c)^2}{a+b+c+3}.\]
Thay $a+b+c = 3$ vào:
\[P \ge \frac{3^2}{3 + 3} = \frac{9}{6} = \frac{3}{2}.\]
Dấu ``='' xảy ra khi và chỉ khi:
\[\frac{a}{b+1} = \frac{b}{c+1} = \frac{c}{a+1} \iff a = b = c = 1.\]
Vậy giá trị nhỏ nhất của biểu thức là $\boxed{\min P = \dfrac{3}{2}}$.
\end{loi_giai_cau}
